\section{Model Class: Does Extraction Require Nonlinearity?}
\label{sec:linvsnonlin}

This section holds the information set fixed and varies the estimator. All
arms are causally tuned by the protocol of Section~\ref{sec:causal_tuning}, so
no arm's hyperparameters ever see the losses they are scored on, and all sit on
the shared panel.

\subsection{The Comparison}
\label{sec:class_table}

\begin{table}[H]
\centering
\caption{Model class at fixed information. Best tree (LightGBM) against best
linear (ridge) on each bucket, causally tuned, shared panel
($n = 218{,}934$). $\Delta_{\text{class}}$ is tree minus linear on the same
feature set. OOS-$R^2$ is raw-variance MSE relative to the shared incumbent;
MZ-$\beta$ is the Mincer--Zarnowitz slope, with $1.0$ denoting an unbiased
forecast.}
\label{tab:class}
\begin{tabular}{lrrrrr}
\toprule
Feature set & LightGBM & ridge & $\Delta_{\text{class}}$ & LGBM OOS-$R^2$ & LGBM MZ-$\beta$ \\
\midrule
all\_features & \textbf{0.12604} & 0.12788 & $-0.00184$ & $+0.059$ & 0.892 \\
moments       & 0.12816 & 0.13085 & $-0.00269$ & $+0.021$ & 0.882 \\
liquidity     & 0.12840 & 0.13120 & $-0.00280$ & $+0.058$ & 0.892 \\
market\_vw    & 0.12963 & 0.13291 & $-0.00328$ & $-0.010$ & 0.862 \\
vol\_demand   & 0.12979 & 0.13337 & $-0.00358$ & $-0.035$ & 0.861 \\
implied\_vol  & 0.13001 & 0.13252 & $-0.00251$ & $-0.034$ & 0.840 \\
sentiment     & 0.13038 & 0.13347 & $-0.00309$ & $-0.063$ & 0.867 \\
market\_ew    & 0.13039 & 0.13319 & $-0.00280$ & $+0.003$ & 0.867 \\
baseline      & 0.13118 & 0.13445 & $-0.00327$ & $-0.040$ & 0.870 \\
\midrule
incumbent (OLS) & \multicolumn{2}{r}{0.13415} & & --- & 0.929 \\
OLS, base-2 ladder & \multicolumn{2}{r}{0.13332} & & & \\
\bottomrule
\end{tabular}
\end{table}

The tree holds a modest edge on every one of the nine feature sets, ranging
from $0.00184$ to $0.00358$, and Section~\ref{sec:pairwise} shows every one of
those nine differences to be statistically resolved under an assumption-free
bound. The consistency of the sign remains the most persuasive part of the
evidence---an exact sign test over the nine sets gives $p = 0.0039$---but the
individual comparisons no longer rest on it alone. The magnitude is small.

\paragraph{The pairwise difference is tested, and it holds.} Every DM statistic
in the battery is computed against the shared incumbent, a third model, so
$t = -15.4$ for LightGBM and $t = -19.9$ for ridge on \texttt{all\_features} do
not by themselves establish that the $0.00184$ between them differs from zero.
Section~\ref{sec:pairwise} supplies the missing test. Because the incumbent
cancels from the pairwise differential, the pairwise standard error is bounded
above by $se_{\mathrm{tree}} + se_{\mathrm{linear}}$ without any assumption
about the dependence between arms, and the resulting bound resolves
\emph{all nine} matched comparisons in the tree's favour after Holm adjustment
(Table~\ref{tab:pairwise_bounds}). The headline pair is the weakest of the
nine, at $|t| \ge 2.19$, $p \le 0.029$; the other eight range from
$|t| \ge 3.73$ to $|t| \ge 5.91$.

Three caveats travel with that result. The reported statistics are bounds, so
the true evidence is stronger by an unmeasured amount. The headline pair is
only just resolved, which means the paper's most-quoted number is its least
secure. And a model confidence set over the 45 arms still requires their
per-bar losses, which are not retained.
\pending{battery MCS once per-bar loss series are persisted}

\subsection{The Pure Nonlinearity Premium, Corrected}
\label{sec:premium}

On the \texttt{baseline} feature set---HAR and calendar only, no exogenous
information---the causally tuned LightGBM reaches $0.13118$ against the
incumbent's $0.13415$, a gap of $0.00297$. This is the quantity earlier drafts
reported as the pure nonlinearity premium, on the grounds that the two arms
share an information set exactly.

Two corrections apply.

\paragraph{Backbone misspecification.} The incumbent is quadrature-limited
(Section~\ref{sec:benchmark_debt}). Measured against the base-2 OLS ladder at
$0.13332$, the premium is
\begin{equation}
  0.13332 - 0.13118 \;=\; 0.00215,
\end{equation}
so approximately 28\% of the headline figure is HAR lag-grid quadrature error
rather than nonlinearity. A tree on a log-spaced boxcar basis can interpolate
between rungs; part of what it was being credited for is doing the linear
model's quadrature for it.

\paragraph{Protocol asymmetry.} The comparison sets a \emph{causally tuned}
LightGBM against an \emph{untuned} OLS. The matched linear arm on the same
features---causally tuned ridge, $0.13445$---is worse than the untuned
incumbent. The premium is therefore measured against whichever linear
specification happens to look best, which is not a clean isolate. We report it
as a range: against the base-2 OLS ladder, $0.00215$; against the causally
tuned ridge on the same features, $0.00327$. The lower figure is the
conservative one and is the number we carry forward.

\subsection{Tuning Is a Non-Lever}
\label{sec:tuning_nonlever}

Three independent measurements agree that hyperparameters are not where the
skill is.

\begin{enumerate}
  \item \textbf{Causal tuning trails fixed defaults on the linear arms.} On the
    \texttt{baseline} feature set, causally tuned ridge ($0.13445$), elastic
    net ($0.13471$) and lasso ($0.13475$) are all \emph{worse} than the fixed
    incumbent at $0.13415$, at DM $t = +5.0$, $+7.5$ and $+7.0$ respectively.
    Selecting a penalty causally, on a leakage-clean validation tail, costs
    more in selection noise than it recovers in fit.
  \item \textbf{An oracle over penalties has almost nothing to give.} A
    test-set oracle---selecting the penalty that minimises realised
    out-of-sample loss, an unattainable upper bound---improves the incumbent by
    at most $0.00015$, an order of magnitude below the information effects.
  \item \textbf{Selection loses to shrinkage at every layer tested.} Heavy
    $\ell_1$ hurts; PCA and PLS truncation hurt; light shrinkage that prunes
    nothing important is the best linear response. This is the estimator-side
    restatement of dense-but-weak.
\end{enumerate}

\subsection{A Cautionary Result: Tuning on the Evaluation Panel}
\label{sec:tuning_leak}

An earlier version of this study selected hyperparameters by minimising QLIKE
on the full walk-forward backtest---the same panel and the same loss used for
evaluation. We report the outcome because it is a useful calibration on the
literature rather than a result about volatility.

Under that protocol the tuned tree ensembles reached QLIKE values around
$0.11$, roughly $0.02$ better than any arm in the causally tuned battery, and
the linear-versus-tree ordering \emph{reversed} relative to the default
configuration. Trial counts differed sharply by model (125 for XGBoost, 30 for
LightGBM, 16 for random forest), so the arms were ranked partly by how many
draws each had taken from the evaluation panel. The apparent effect,
$\approx 0.03$ QLIKE, is an order of magnitude larger than every genuine
model-class effect in Table~\ref{tab:class}.

The introduction claims that published disagreements about whether nonlinear
models help on this task are partly mechanical. This is the mechanism,
measured on our own panel: a tuning protocol that touches the evaluation set
can manufacture a model-class reversal several times larger than the effect
under study. We regard the earlier numbers as uninterpretable and they appear
nowhere else in this paper.

\subsection{The Metrics Disagree}
\label{sec:metric_disagreement}

QLIKE is our primary ranker, on the standard grounds that it is robust to
noise in the volatility proxy \citep{Patton2011}. That choice is consequential
here in a way that a single reported metric would hide.

Across the 45 causally tuned arms, the rank correlation between QLIKE and
raw-variance OOS-$R^2$ is $+0.152$ ($p = 0.32$). Since lower QLIKE and higher
$R^2$ both denote better forecasts, agreement would produce a \emph{negative}
correlation. The observed value is near zero, with the wrong sign, and is not
distinguishable from no relationship at all: the two metrics do not rank this
battery in a related way.

Concretely, \textbf{14 of the 42 arms that improve on the incumbent under QLIKE
are worse than the incumbent under MSE.} These include the second-best arm
overall (\texttt{xgb / all\_features}, QLIKE $0.12641$, OOS-$R^2 = -0.062$) and
the best linear arm (\texttt{ridge / all\_features}, QLIKE $0.12788$,
OOS-$R^2 = -0.058$). Only the single best arm, \texttt{lgbm / all\_features},
improves on both.

Raw-variance MSE on this panel is dominated by a small number of extreme bars,
most of them in March 2020, which is the standard argument for preferring
QLIKE. We accept that argument but note that it is a \emph{post hoc} defence of
the metric under which our own ordering looks best, and that it should be
settled with evidence rather than assertion.
\pending{MSE recomputed with the COVID window excluded, reported alongside the
full-sample MSE, so that the reader can see whether the QLIKE/MSE divergence is
a tail-domination artifact or a genuine disagreement about central performance}

\subsection{QLIKE Rewards Attenuation}
\label{sec:calibration}

Every arm in the battery is miscalibrated in the same direction: Mincer--Zarnowitz
slopes run from $0.840$ to $0.941$, all below unity, indicating forecasts that
are systematically too flat relative to realised variance.

The pattern within that is the informative part. Across the 45 causally tuned
arms, the rank correlation between QLIKE and MZ-$\beta$ is $\mathbf{+0.741}$
($p = 5.8 \times 10^{-9}$):
\emph{the arms that score best on QLIKE are the most attenuated}. The best arm,
\texttt{lgbm / all\_features}, has $\beta = 0.892$ against the incumbent's
$0.929$; the second-best, \texttt{xgb / all\_features}, has the worst
calibration in the entire battery at $\beta = 0.840$.

This does not invalidate the ranking---QLIKE is a strictly consistent scoring
rule and the ordering it induces is the one it is designed to induce---but it
does change what the ranking means. QLIKE penalises under-prediction far more
heavily than over-prediction, so a forecaster that shades toward the
conditional mean of the loss rather than the conditional mean of the target is
rewarded. A practitioner who cares about calibrated variance forecasts, rather
than about QLIKE, would rank these arms differently, and the incumbent would
move up. We report this as a finding about the loss function rather than about
the models.

\subsection{The Nonlinearity Saturates}
\label{sec:saturation}

The recoverable nonlinear structure has a ceiling, and the ceiling is low.

\begin{itemize}
  \item Restricting the tree to additive-plus-pairwise structure anchored to
    the session clock recovers essentially all of the tree's advantage;
    deeper interaction capacity does not pay.
  \item Deep sequence models and mixture-of-experts routing fail to improve on
    the saturated form.
  \item On the kernel dictionary of Section~\ref{sec:descriptive}, trees fail
    as a \emph{direct} forecaster, and fail decisively: per-bar DM on the
    campaign tile puts the extrapolation clamp at $+0.077$ for LightGBM
    ($t = 6.46$) and $+0.102$ for XGBoost ($t = 6.96$), the largest single
    effect in this paper (Section~\ref{sec:tile_inference}). As a
    \emph{residual reader} the tree is worse than a smooth linear reader by
    $\approx 0.003$, but that difference is \emph{not} statistically resolved
    ($t = 1.06$ and $1.11$). We therefore claim the failure for the direct seat
    only. The nonlinearity that remains is smooth and diffuse rather than
    tree-shaped, which is consistent with---though not identified by---a
    second-order feedback.
\end{itemize}

Under a generic machine-learning framing this reads as an unexplained ceiling.
Under the backbone of Section~\ref{sec:hawkes} it reads as a prediction. The
quadratic Hawkes model of \citet{BlancDonierBouchaud2017} makes the feedback a
quadratic form in the past rather than a linear functional---generating
clustering, fat tails and the Zumbach effect endogenously---and what it implies
about extractable structure is exactly additive plus pairwise and no more. The
saturation we measure is the saturation that model predicts.

We are careful about the strength of this claim. We have not fitted a quadratic
Hawkes model, and the agreement is qualitative: a measured ceiling at pairwise
order matching a model class whose feedback is of second order. It is
consistency, not identification. \pending{fitting the quadratic feedback form
directly and testing it against the saturated tree would convert this from
consistency to evidence}

\subsection{Summary}
\label{sec:class_summary}

Within the dense-but-weak regime, nonlinearity buys a small, consistent,
bounded amount---between $0.00184$ and $0.00358$ depending on feature set, with
a pure premium of $0.00215$ against a correctly specified linear backbone---and
hyperparameters buy nothing. Set against the joint information effect of
$0.00627$, the ordering
\emph{information} $>$ \emph{model class} $>$ \emph{hyperparameters} holds on
this panel by roughly a factor of three at each step. That ordering is the
paper's central practical claim, subject to the pairwise significance test that
Section~\ref{sec:class_table} flags as outstanding.
